\begin{helper}[Last-return data contradict global maximality]
For every defective color $\delta$, no
\noderef{ScheideDefectiveReturnData} for $\delta$ exists.
\end{helper}
\begin{proof}
Suppose that a return datum $R$ exists.  Put
$X=\noderef{TreeSequenceVertexSet}(T)$,
$i=\gammaIndex(\delta)$, $j=R.j$,
$\gamma=\gamma_i$, and $\eta=\gamma_j$.  Let
$\varphi_2$ be the first final coloring, obtained from the
$\alpha _0,\gamma$ \noderef{OutsideTreeColorSwap} followed by the
\noderef{TwoColorComponentSwap} through $v^0_\delta$, and let $\psi_2$
be the second final coloring from
\noderef{ScheideDefectiveSecondTerminalComponent}.  The color $\alpha _0$,
the vertices $v^0_\delta,u^0_\delta$, and the edge $f^0_\delta$ are the
fields of \noderef{ScheideDefectiveSetup}; the index $j$, vertices
$R.v_1,R.u_1$, edge $R.f_1$, endpoint equations, and color distinctions
used below are the fields of \noderef{ScheideDefectiveReturnData}.

Choose the partial Tashkinov trees $S_\gamma,S_\eta$ from
\noderef{ScheideDefectiveDeletedTreePair}, and choose
$f_2=z_2u_2$ from \noderef{ScheideDefectiveThirdBoundaryEdge}.  Thus
\[
 V(S_\gamma)=(X\setminus U_\gamma)\cup\{R.u_1\},qquad
 V(S_\eta)=(X\setminus U_\eta)\cup\{u^0_\delta\},             \tag{1}
\]
$u_2\notin X$, $u_2\ne R.u_1,u^0_\delta$, and $f_2$ has color
$\delta$ in both final colorings.

Append $(f_2,u_2)$ to both sequences.  In the first sequence, $z_2$ is
an earlier vertex because it belongs to
$X\setminus U_\gamma$; in the second it belongs to
$X\setminus U_\eta$: this is exactly the assertion $z_2\in Z$ from
\noderef{ScheideDefectiveThirdBoundaryEdge}.  The endpoint $u_2$ belongs
to neither partial vertex set in (1), since $u_2\notin X$ and
$u_2\ne R.u_1,u^0_\delta$.  Hence the exact endpoint equation
$\operatorname{ends}(f_2)=\{z_2,u_2\}$ puts $f_2$ in
\noderef{BoundaryEdges}$(G,V(S_\gamma))$ and in
$\noderef{BoundaryEdges}(G,V(S_\eta))$.  Therefore
\noderef{TreeSequenceBoundaryEdgeAppendExists} itself supplies a fresh
exact append to each partial tree and produces
tree-sequences
\[
 \widehat T_\gamma,\qquad \widehat T_\eta
\]
with vertex sets
\[
 V(\widehat T_\gamma)
   =(X\setminus U_\gamma)\cup\{R.u_1,u_2\},qquad
 V(\widehat T_\eta)
   =(X\setminus U_\eta)\cup\{u^0_\delta,u_2\}.                \tag{2}
\]
These equalities are the vertex-set clause of
\noderef{TreeSequenceAdjoinEdge} applied to (1); its
\noderef{TreeSequenceInitialSegment} clause retains every old displayed
identity, while the boundary-edge construction supplies the new identity.

We check the new Tashkinov clauses rather than infer them from
cardinality.  For the first recoloring,
\noderef{ScheideDefectiveTerminalComponent} identifies the exact switched
component.  The setup \noderef{ScheideDefectiveSetup} says that $\gamma$
is missing at $v^0_\delta$; the outside swap fixes every incident edge of
this tree vertex by \noderef{OutsideTreeColorSwap}; and
\noderef{TwoColorComponentSwapMissingBehavior} then says that $\delta$ is
missing there after the component switch.  The exit datum in the setup and
\noderef{ScheideStructuralExitsInBase} put $v^0_\delta$ in the retained
base, so it occurs earlier than the appended edge.  For the second
recoloring,
\noderef{ScheideDefectiveSecondTerminalComponent} says that $\delta$ is
missing at $R.v_1$, and \noderef{ScheideDefectiveReturnData} puts this
vertex in the retained base.  In both cases $f_2$ has color $\delta$, so
the corresponding retained vertex is an earlier missing witness.  Every
old Tashkinov clause comes from the two partial trees and is carried to the
append by the initial-segment maps.  This verifies the unique new-edge
clause, and hence every clause of \noderef{TashkinovTree}:
\[
 \noderef{TashkinovTree}(G,e_0,\varphi_2,\widehat T_\gamma),
 \qquad
 \noderef{TashkinovTree}(G,e_0,\psi_2,\widehat T_\eta).       \tag{3}
\]

The exact position fields of \noderef{ScheideStructuralConfiguration}
and vertex injectivity in \noderef{TreeSequence} show that each deleted
block is a two-element subset of $X$.  In the first formula in (2),
$R.u_1,u_2$ are distinct and outside $X$; in the second,
$u^0_\delta,u_2$ are distinct and outside $X$.  These are respectively
the outside fields of \noderef{ScheideDefectiveReturnData} and
\noderef{ScheideDefectiveSetup}, and the exclusions in
\noderef{ScheideDefectiveThirdBoundaryEdge}.  Thus both trees in (2) have
exactly $|X|$ vertices.

We now rebuild their global maximality field by field.  For $\varphi_2$,
the first outside swap is proper and leaves exactly $e_0$ uncolored by
\noderef{ScheideDefectiveTerminalComponent}; its two component colors are
distinct by the return datum, so
\noderef{TwoColorComponentSwapProper} gives properness after the switch.
The definition of \noderef{TwoColorComponentSwap} preserves whether each
exact identity is colored or uncolored, so the definition of
\noderef{IsOneEdgeUncolored} still singles out exactly $e_0$.  The
corresponding properness and exactly-one-uncolored facts for $\psi_2$ are
stated directly in \noderef{ScheideDefectiveSecondTerminalComponent}.
Thus in each case the definition of
\noderef{IsProperEdgeColoring}, together with the unchanged criticality,
chromatic-index, and maximum-degree fields of the original
\noderef{KTriple}, gives a $K$-triple for the new coloring.
Equations (2)--(3) and the original equality
$|X|=\noderef{TashkinovOrder}(G)$ then satisfy every clause of
\noderef{GlobalMaximumTriple}.  Thus both new triples are globally
maximum.  Applying \noderef{ScheideGlobalMaximumOrder} to them makes
the two sets in (2) \noderef{ColoredVertexSetElementary} and
\noderef{ColoredVertexSetClosed} in their respective colorings.

It remains to contradict the simultaneous global maximality.  The exact
base-position equations in \noderef{ScheideStructuralConfiguration} and
the attachment law in \noderef{TreeSequence} identify the endpoints of
the first edge $e_0$ as the distinct first two base vertices $y_0,y_1$;
in particular each is \noderef{Incident} with $e_0$.
Apply \noderef{ScheideGlobalMaximumMissingCounts} to the original global
maximum triple at each endpoint.  Each missing-color set has at least
two members.  The original set $X$ is elementary by
\noderef{ScheideGlobalMaximumOrder}; by the definition of
\noderef{ColoredVertexSetElementary}, the two sets are disjoint, so their
union has at least four colors.  Choose a color $\beta$ in this union
outside $\{\alpha _0,\gamma,\eta\}$, and let
$y_\beta\in\{y_0,y_1\}$ miss it.  The defective color $\delta$ is
represented by a member of \noderef{DefectiveColors}.  Unfolding
\noderef{DefectiveColor} and \noderef{BoundaryColorEdges}, it therefore
has a $\delta$-colored boundary identity of $X$.  If $\delta$ were
missing at $y_\beta$, that identity would violate the definition of
\noderef{ColoredVertexSetClosed} for the original tree.  Thus $\beta$
cannot equal $\delta$, and
\[
 \beta\notin\{\alpha _0,\gamma,\eta,\delta\}.                \tag{4}
\]

The definition of \noderef{OutsideTreeColorSwap} and, at the second
stage, that of \noderef{TwoColorComponentSwap}, shows edge by edge that
each operation producing either final coloring transposes two colors from
the set on the right of (4).  Hence every exact edge has color $\beta$ in
a final coloring exactly when it had color $\beta$ originally;
$\beta$ remains missing at the retained base vertex $y_\beta$ in both
new maximum trees.  By \noderef{ColoredVertexSetElementary}, the
distinct vertices $y_\beta,u_2$ cannot both miss $\beta$ in the first
new tree.  Thus $\beta$ is not missing at $u_2$; by the complementary
definitions of \noderef{MissingColors} and \noderef{PresentColors}, it is
present there.  Choose its exact incident $\beta$-edge $b$.
If its other endpoint lay outside the first new tree, then the endpoint
pair would put $b$ in \noderef{BoundaryEdges} of that tree, contradicting
\noderef{ColoredVertexSetClosed} with the missing witness $y_\beta$.
Thus the other endpoint lies in the first new tree.  It cannot lie in
$X$: since $u_2\notin X$, its endpoint pair would instead put $b$ in the
boundary of the original closed set $X$, where $\beta$ is missing.
From the first equality in (2), and
\noderef{LooplessMultigraph},
the other endpoint is therefore exactly $R.u_1$:
\[
 \operatorname{ends}(b)=\{u_2,R.u_1\}.                       \tag{5}
\]

Equation (4) shows that the same identity $b$ has color $\beta$ under
$\psi_2$.  The vertex $R.u_1$ is absent from the second set in (2): it
lies outside $X$, it differs from $u_2$ by the third-edge helper, and it
differs from $u^0_\delta$ as follows.  First $R.f_1\ne f^0_\delta$:
equality, together with the exact endpoint pairs in
\noderef{ScheideDefectiveReturnData} and
\noderef{ScheideDefectiveSetup}, would either identify the distinct tree
vertices $R.v_1,v^0_\delta$ or identify a vertex in $X$ with one outside
$X$.  If nevertheless $R.u_1=u^0_\delta$, the two distinct identities
$R.f_1,f^0_\delta$ would be \noderef{AdjacentEdges}, since both are
\noderef{Incident} with this common outside endpoint.  Both have color
$\delta$ under the original coloring, contradicting
\noderef{IsProperEdgeColoring} for the original $K$-triple.  Therefore
$R.u_1\ne u^0_\delta$, as required.

Now (5), the membership of $u_2$ in the second tree, and the just-proved
nonmembership of $R.u_1$ put $b$ directly in
\noderef{BoundaryEdges} of the second tree.  Hence
\noderef{TreeSequenceBoundaryEdgeAppendExists} supplies a fresh exact
append by $(b,R.u_1)$.  Its initial-segment clauses preserve all old
Tashkinov witnesses, and the new $\beta$-edge has the retained earlier
missing witness $y_\beta$.  The result is therefore a
\noderef{TashkinovTree} for the same edge and coloring.

Its vertex set strictly contains the second globally maximum tree's
vertex set because it adds $R.u_1$.  This contradicts
\noderef{ScheideGlobalMaximumFixedPairMaximal}, equivalently the maximum
order bound encoded by \noderef{TashkinovOrder}.  Therefore the assumed
return datum cannot exist.
\end{proof}
